\begin{frame}{이번 주는 어떻게 진행되는가}
  \small
  \begin{tabular}{@{}clp{0.6\textwidth}@{}}
    \toprule
    순서 & 내용 & 시간 배분 \\
    \midrule
    1 & 다른 두 팀의 발표 청취(팀당 5$\sim$7분, 8.1 형식 그대로) +
      보고서·슬라이드 열람 & 약 30분 \\
    2 & \textbf{작성 리뷰 2통} --- 팀당 체크리스트 4항목 + \textbf{반박
      가능한 질문 1개 이상}(개념의 장·절 명시) & 발표 직후(각 10분) \\
    3 & 질의·토론 --- 리뷰어가 쓴 질문에 발표팀이 그 자리에서 답변 &
      나머지 시간 \\
    \bottomrule
  \end{tabular}
  \vskip0.8em
  \begin{itemize}
    \item 핵심은 2번의 \textbf{작성 리뷰}. 3번 토론은 ``발표팀이 그 질문에
      \textbf{실제로} 답할 수 있는가''를 가르는 자리.
    \item \kb{3점 질문} = 답하려면 \textbf{새로운 계산·실험}(시드 재실행,
      $\varepsilon=0$ 재평가, $\alpha$ 민감성 등)을 꺼내야 하는 질문.
      ``네, 알겠습니다''로 넘어가는 질문은 2점에 머물러.
    \item 8.1에서 ``수식적 정당화''가 발표팀의 \emph{의무}였다면, 이번 주
      그 정당화가 \emph{타당한지}를 리뷰어가 \emph{심사}.
  \end{itemize}
\end{frame}
